% GENERATED FILE. Edit canonical lessons, then run npm run generate:books.
% canonical-source-hash: fnv1a64:fc2c8c7d82c3b31b
% canonical-lessons: LA-C458-angiportum, LA-C458-vicinia, LA-C458-latrina, LA-C458-cloaca, LA-C458-castellum

\chapter{Alley, Neighbourhood, Privy, Sewer, Water tower}
\label{ch:la-alley-neighbourhood-privy-sewer-water-tower}
\hlchaptermodality{458}

\begin{chapteropening}
\textbf{By the end of this chapter:} \emph{I can talk about the back streets and services of a town: an alley, the neighbourhood, a privy, the sewer and a water tower.}

Complete the last lesson of chapter 458: I can talk about the back streets and services of a town: an alley, the neighbourhood, a privy, the sewer and a water tower.
\end{chapteropening}

\section[angiportum, angiportī]{angiportum, angiportī — an alley, a narrow lane}

\label{lesson:LA-C458-angiportum}



\begin{quote}
Before the new one: say the Latin for a timber, then the Latin for a beam.
\end{quote}

\subsection*{You'll want to know: angiportum, angiportī}
\textbf{angiportum, angiportī} — \textquotedblleft{}an alley, a narrow lane\textquotedblright{}.

\begin{quote}
Nocte angiportum obscūrum et perīculōsum est.
\end{quote}

\begin{quote}
\emph{At night the alley is dark and dangerous.}
\end{quote}

Second declension, neuter: \emph{angiportum, angiportī n.}; also \emph{angiportus, angiportūs m.} A narrow side street between houses.

\begin{cousinweb}
\emph{angiportum} is built, as Lewis \& Short explain under \emph{angiportus}, on the root of \emph{angustus} 'narrow' and on \emph{portus} in its old sense 'passage'. EtymDB does not list it.
\end{cousinweb}

\begin{grammarlens}[title={What you've built}]
One more word for telling, explaining and describing at B1.
\end{grammarlens}

\subsection*{Your turn}
\begin{itemize}
\raggedright
  \item \emph{Say it:} \emph{angiportum, angiportī}
  \item \emph{Say it:} \emph{angiportum, angiportī}, once more
  \item \emph{Say it:} a sentence of your own with \emph{angiportum, angiportī}
  \item \emph{Recall:} say the Latin for a timber, then the Latin for a beam, then say \emph{angiportum, angiportī} again
\end{itemize}

\begin{tcolorbox}[breakable,colback=teal!4,colframe=teal!35!black,title={Before you move on}]
What does angiportum, angiportī mean? (\textquotedblleft{}An alley, a narrow lane\textquotedblright{}.) Say it once more.
\end{tcolorbox}
\section[vīcīnia, vīcīniae]{vīcīnia, vīcīniae — a neighbourhood; the neighbours}

\label{lesson:LA-C458-vicinia}



\begin{quote}
Before the new one: say the Latin for a beam, then the Latin for an alley.
\end{quote}

\subsection*{You'll want to know: vīcīnia, vīcīniae}
\textbf{vīcīnia, vīcīniae} — \textquotedblleft{}a neighbourhood; the neighbours\textquotedblright{}.

\begin{quote}
In nostrā vīcīniā taberna bona est.
\end{quote}

\begin{quote}
\emph{In our neighbourhood there is a good shop.}
\end{quote}

First declension, feminine: \emph{vīcīnia, vīcīniae f.}, from \emph{vīcīnus} 'neighbour'. English \emph{vicinity} comes from the related \emph{vīcīnitās}.

\begin{cousinweb}
\emph{vīcīnia} is built on \emph{vīcīnus} 'neighbour' with the suffix \emph{-ia} (EtymDB; Lewis \& Short), and \emph{vīcīnus} comes from \emph{vīcus} 'street, village'.
\end{cousinweb}

\begin{grammarlens}[title={What you've built}]
One more word for telling, explaining and describing at B1.
\end{grammarlens}

\subsection*{Your turn}
\begin{itemize}
\raggedright
  \item \emph{Say it:} \emph{vīcīnia, vīcīniae}
  \item \emph{Say it:} \emph{vīcīnia, vīcīniae}, once more
  \item \emph{Say it:} a sentence of your own with \emph{vīcīnia, vīcīniae}
  \item \emph{Recall:} say the Latin for a beam, then the Latin for an alley, then say \emph{vīcīnia, vīcīniae} again
\end{itemize}

\begin{tcolorbox}[breakable,colback=teal!4,colframe=teal!35!black,title={Before you move on}]
What does vīcīnia, vīcīniae mean? (\textquotedblleft{}A neighbourhood; the neighbours\textquotedblright{}.) Say it once more.
\end{tcolorbox}
\section[lātrīna, lātrīnae]{lātrīna, lātrīnae — a latrine, a privy, a lavatory}

\label{lesson:LA-C458-latrina}



\begin{quote}
Before the new one: say the Latin for an alley, then the Latin for a neighbourhood.
\end{quote}

\subsection*{You'll want to know: lātrīna, lātrīnae}
\textbf{lātrīna, lātrīnae} — \textquotedblleft{}a latrine, a privy, a lavatory\textquotedblright{}.

\begin{quote}
In thermīs lātrīna pūblica est.
\end{quote}

\begin{quote}
\emph{In the baths there is a public lavatory.}
\end{quote}

First declension, feminine: \emph{lātrīna, lātrīnae f.} Lewis \& Short note that it first meant a bath. English \emph{latrine} comes from it.

\begin{cousinweb}
\emph{lātrīna} is shortened from \emph{lavātrīna} 'a place for washing', built on \emph{lavō} 'to wash' (Lewis \& Short). English \emph{latrine} comes from it.
\end{cousinweb}

\begin{grammarlens}[title={What you've built}]
One more word for telling, explaining and describing at B1.
\end{grammarlens}

\subsection*{Your turn}
\begin{itemize}
\raggedright
  \item \emph{Say it:} \emph{lātrīna, lātrīnae}
  \item \emph{Say it:} \emph{lātrīna, lātrīnae}, once more
  \item \emph{Say it:} a sentence of your own with \emph{lātrīna, lātrīnae}
  \item \emph{Recall:} say the Latin for an alley, then the Latin for a neighbourhood, then say \emph{lātrīna, lātrīnae} again
\end{itemize}

\begin{tcolorbox}[breakable,colback=teal!4,colframe=teal!35!black,title={Before you move on}]
What does lātrīna, lātrīnae mean? (\textquotedblleft{}A latrine, a privy, a lavatory\textquotedblright{}.) Say it once more.
\end{tcolorbox}
\section[cloāca, cloācae]{cloāca, cloācae — a sewer, a drain}

\label{lesson:LA-C458-cloaca}



\begin{quote}
Before the new one: say the Latin for a neighbourhood, then the Latin for a latrine.
\end{quote}

\subsection*{You'll want to know: cloāca, cloācae}
\textbf{cloāca, cloācae} — \textquotedblleft{}a sewer, a drain\textquotedblright{}.

\begin{quote}
Aqua sordida per cloācam in flūmen it.
\end{quote}

\begin{quote}
\emph{Dirty water runs through the sewer into the river.}
\end{quote}

First declension, feminine: \emph{cloāca, cloācae f.} The great sewers of old towns still carry water.

\begin{cousinweb}
\emph{cloāca} is built on an old verb \emph{cluō} 'to cleanse, to purge' (EtymDB; Lewis \& Short), so it is a channel that washes the city clean. English keeps \emph{cloaca} as a term in zoology.
\end{cousinweb}

\begin{grammarlens}[title={What you've built}]
One more word for telling, explaining and describing at B1.
\end{grammarlens}

\subsection*{Your turn}
\begin{itemize}
\raggedright
  \item \emph{Say it:} \emph{cloāca, cloācae}
  \item \emph{Say it:} \emph{cloāca, cloācae}, once more
  \item \emph{Say it:} a sentence of your own with \emph{cloāca, cloācae}
  \item \emph{Recall:} say the Latin for a neighbourhood, then the Latin for a latrine, then say \emph{cloāca, cloācae} again
\end{itemize}

\begin{tcolorbox}[breakable,colback=teal!4,colframe=teal!35!black,title={Before you move on}]
What does cloāca, cloācae mean? (\textquotedblleft{}A sewer, a drain\textquotedblright{}.) Say it once more.
\end{tcolorbox}
\section[castellum, castellī]{castellum, castellī — a water tower, a reservoir; a fort, a stronghold}

\label{lesson:LA-C458-castellum}



\begin{quote}
Before the new one: say the Latin for a latrine, then the Latin for a sewer.
\end{quote}

\subsection*{You'll want to know: castellum, castellī}
\textbf{castellum, castellī} — \textquotedblleft{}a water tower, a reservoir; a fort, a stronghold\textquotedblright{}.

\begin{quote}
Aqua ex aquaeductū in castellum venit et inde in urbem.
\end{quote}

\begin{quote}
\emph{Water comes from the aqueduct into the water tower and from there into the city.}
\end{quote}

Second declension, neuter: \emph{castellum, castellī n.}, a diminutive of \emph{castrum} 'fort'. The \emph{castellum aquae} shared the water out to the pipes. English \emph{castle} comes from it.

\begin{cousinweb}
\emph{castellum} is a diminutive of \emph{castrum} 'fort, fortified place' (EtymDB; Lewis \& Short), and Lewis \& Short also give the water reservoir of an aqueduct. English \emph{castle} comes from it.
\end{cousinweb}

\begin{grammarlens}[title={What you've built}]
One more word for telling, explaining and describing at B1.
\end{grammarlens}

\subsection*{Your turn}
\begin{itemize}
\raggedright
  \item \emph{Say it:} \emph{castellum, castellī}
  \item \emph{Say it:} \emph{castellum, castellī}, once more
  \item \emph{Say it:} a sentence of your own with \emph{castellum, castellī}
  \item \emph{Recall:} say the Latin for a latrine, then the Latin for a sewer, then say \emph{castellum, castellī} again
\end{itemize}

\begin{tcolorbox}[breakable,colback=teal!4,colframe=teal!35!black,title={Before you move on}]
What does castellum, castellī mean? (\textquotedblleft{}A water tower, a reservoir; a fort, a stronghold\textquotedblright{}.) Say it once more.
\end{tcolorbox}
